%% ma: L11-B15-0095
%% lop: 11
%% chuong: 5
%% bai: 15
%% dang: 11.15.5
%% loai: TLN
%% muc: VDC
%% dap_an: "3"
%% nguon: "tai-lieu-hoc-tap-mon-toan-lop-11-theo-chuong-trinh-giao-duc-pho-thong-2018.pdf (Lớp Toán Thầy Duy, Phạm Lê Duy, THPT Chu Văn An), Chương 3 Bài 1, Phần 3, Câu 17 (tr.346)"
%% kien_thuc: "Lớp 11 Bài 15 (giới hạn của dãy số)"
%% trang_thai: cho-duyet
%% ly_do: "Dạng mới đề xuất (11.15.x), chờ giáo viên duyệt dạng; tài liệu gốc không có lời giải mẫu, đáp án do kho tự giải và kiểm bằng sympy."
\begin{ex}
Biết giới hạn \[\lim\dfrac{\sqrt[3]{2n^2-n^3}+n}{\sqrt{n^2+n}-n}=\lim\dfrac{2\left(\sqrt{1+\dfrac1n}+1\right)}{\sqrt[3]{\left(\dfrac an-b\right)^2}+1-\sqrt[3]{\dfrac2n-1}}\] với $a;b$ là các số nguyên dương. Tính $S=a+b$.
\shortans{3}
\loigiai{\sloppy Đặt $x=\dfrac2n-1$, khi đó $\sqrt[3]{2n^2-n^3}=n\sqrt[3]x$ và \[\sqrt[3]{2n^2-n^3}+n=n\left(\sqrt[3]x+1\right)=\dfrac{n(x+1)}{\sqrt[3]{x^2}-\sqrt[3]x+1}=\dfrac{2}{\sqrt[3]{x^2}-\sqrt[3]x+1}.\] Mẫu: $\sqrt{n^2+n}-n=\dfrac1{\sqrt{1+\frac1n}+1}$. Vậy phân thức bằng \[\dfrac{2\left(\sqrt{1+\frac1n}+1\right)}{\sqrt[3]{\left(\frac2n-1\right)^2}+1-\sqrt[3]{\frac2n-1}},\] tức $a=2$, $b=1$ và $S=3$.}
\end{ex}
